\section{Assembling a self-dual matroid}\label{sec:assembly}

We now choose the prototype classes to which the local rules will be applied.
There are two requirements: every resulting basis must impose the desired
strict inequality on ordinal ranks, and enough classes must be present to
allow maximal independent extensions inside every subset of $E_0$.
The interval criterion used for the second requirement is due to
Bowler and Geschke~\cite[Theorem~6]{BowlerGeschke2016}; we include the full
verification of the matroid axioms after constructing the bases.

\subsection{Admissible sets and two monotone zones}

For $X\subseteq E_0=L\dot\cup R$, put
\[
 u(X)=X_L,\qquad y(X)=D\setminus X_R.
\]
Call $X$ \emph{admissible} if either
\begin{equation}\label{eq:admissible}
 \begin{aligned}
 &u(X),y(X)\in\mathcal U
       &&\text{and }\rho(u(X))>\rho(y(X)),\\
 \text{or }\quad
 &u(X),y(X)\notin\mathcal U
       &&\text{and }\rho(D\setminus u(X))>
                         \rho(D\setminus y(X)).
 \end{aligned}
\end{equation}
All ranks appearing here are evaluated on members of $\mathcal U$.
Define the \emph{Upper zone} to consist of those $X$ for which $u(X)$ is
large and either $y(X)$ is small, or $y(X)$ is large and
$\rho(u(X))\le\rho(y(X))$.
Define the \emph{Lower zone} by requiring $X^c$ to lie in the Upper zone.
Complements of subsets of $E_0$ are taken in $E_0$.

\begin{lemma}\label{lem:zones}
The Upper zone, Lower zone, and admissible sets partition $2^{E_0}$.
The Upper zone is upward closed, the Lower zone is downward closed,
and all three statuses are invariant under symmetric differences in
$\mathcal K_0$.
Complementation exchanges the two zones and preserves admissibility.
Moreover, $\varnothing$ is Lower, $E_0$ is Upper, and every admissible
set and its complement are $\mathcal J_0$-positive.
\end{lemma}

\begin{proof}
Write $u=u(X)$ and $y=y(X)$.
When $u$ is large and $y$ small, $X$ is Upper; when $u$ is small and $y$
large, it is Lower.
If both are large, comparison of their ranks partitions the possibilities
into Upper and admissible.
If both are small, comparison of the ranks of their complements partitions
the possibilities into Lower and admissible.
Complementing $X$ replaces both $u$ and $y$ by their complements, proving
the asserted symmetries.

Suppose that $X\subseteq X'$ and $X$ is Upper.
Then $u\subseteq u'$ and $y'\subseteq y$, with $u'$ large.
If $y'$ is small, $X'$ is Upper by definition.
Otherwise $y$ is large as well, and Lemma~\ref{lem:rank} gives
\[
 \rho(u')\le\rho(u)\le\rho(y)\le\rho(y'),
\]
so again $X'$ is Upper.
Downward closure of the Lower zone follows by complementation.
The invariance assertion follows from Lemma~\ref{lem:rank} in the two
coordinates.
The statuses of $\varnothing$ and $E_0$ follow directly from the definitions.

If the first line of~\eqref{eq:admissible} holds, $X_L=u$ and
$(X^c)_R=y$ are large, hence $\mathcal J$-positive.
If the second line holds, $X_R=D\setminus y$ and
$(X^c)_L=D\setminus u$ are large instead.
In either case $X$ and $X^c$ are $\mathcal J_0$-positive.
\end{proof}

\begin{lemma}[Crossing the zones]\label{lem:zone-crossing}
If $A\subseteq Z\subseteq E_0$, with $A$ Lower and $Z$ Upper, then
there is an admissible $B$ satisfying $A\subseteq B\subseteq Z$.
\end{lemma}

\begin{proof}
The sets $s=A_L$ and $X=Z_L$ are respectively small and large, and
$s\subseteq X$.
We keep $B_R=A_R$, so write $y=D\setminus A_R$.
If $y$ is large, Lemma~\ref{lem:thinning} gives a large $u$ between
$s$ and $X$ with $\rho(u)>\rho(y)$; take $B_L=u$.
If $y$ is small, apply the same Lemma between the small set
$D\setminus X$ and the large set $D\setminus s$ to obtain a large $v$
with $\rho(v)>\rho(D\setminus y)$.
Taking $B_L=D\setminus v$ gives the second line
of~\eqref{eq:admissible}.
Both constructions respect the two endpoint inclusions.
\end{proof}

\subsection{Reserving separation from earlier classes}

Different prototypes must have positive differences in both directions.
This will ensure that their local bases cannot contain one another.
The next Lemma arranges these two inequalities while preserving the
endpoint statuses needed in Lemma~\ref{lem:zone-crossing}.

\begin{lemma}\label{lem:insert-prototype}
Let $\mathcal T$ be a family of fewer than $\mathfrak c$ admissible subsets
of $E_0$.
Suppose $A\subseteq Z\subseteq E_0$, that $A$ is not Upper and $Z$ is
not Lower, and that, for every $T\in\mathcal T$,
\begin{equation}\label{eq:available-separation}
 T\setminus A\notin\mathcal J_0,
 \qquad Z\setminus T\notin\mathcal J_0.
\end{equation}
Then there is an admissible $B$ between $A$ and $Z$ such that
$B\setminus T$ and $T\setminus B$ are $\mathcal J_0$-positive for every
$T\in\mathcal T$.
\end{lemma}

\begin{proof}
Put $G=Z\setminus A$.
For every $T\in\mathcal T$ with $A\setminus T\in\mathcal J_0$,
include the trace $G\setminus T$ in a list.
It is positive by~\eqref{eq:available-separation}, since
$Z\setminus T=(A\setminus T)\cup(G\setminus T)$.
For every $T$ with $T\setminus Z\in\mathcal J_0$, also include the
positive trace $T\cap G$, using
$T\setminus A=(T\setminus Z)\cup(T\cap G)$.
There are fewer than $\mathfrak c$ traces.
Lemma~\ref{lem:reservations} supplies disjoint $\mathcal K_0$-small sets
$H_0,H_1$ that meet every listed trace $\mathcal J_0$-positively.
An empty list presents no restriction on their choice.

Modify only the gap:
\begin{equation}\label{eq:gap-reservations}
 A'=A\cup(H_0\cap G),\qquad
 Z'=Z\setminus(H_1\cap G).
\end{equation}
Because $H_0$ and $H_1$ are disjoint and $G$ is disjoint from $A$, we have
$A'\subseteq Z'$.
The endpoint changes are $\mathcal K_0$-small, so $A'$ is not Upper and
$Z'$ is not Lower by Lemma~\ref{lem:zones}.
If an endpoint is admissible, take it for $B$.
Otherwise $A'$ is Lower and $Z'$ Upper, and
Lemma~\ref{lem:zone-crossing} gives an admissible $B$ between them.

For any $T\in\mathcal T$, if $A\setminus T$ is positive then so is
$B\setminus T$.
Otherwise the reserved positive set $H_0\cap(G\setminus T)$ lies in
$B\setminus T$.
Likewise, if $T\setminus Z$ is positive it is already contained in
$T\setminus B$; otherwise the reserved positive set $H_1\cap T\cap G$
is contained in $T\setminus B$.
This proves both inequalities.
\end{proof}

\subsection{Meeting every interval}

We choose admissible prototypes in complementary pairs.
For each first prototype $T$ we retain the local bases
$\mathcal B^-(T)$, and for $T^c$ we retain $\mathcal B^+(T^c)$.
These two families are complements of one another by the local rules.
Once a pair has been inserted, its rule and all its bases remain fixed.
The aim is that, for every interval $I\subseteq Y\subseteq E_0$ with infinite gap,
some basis $B$ satisfies the following alternative:
\begin{equation}\label{eq:interval-alternative}
 B\subseteq I,\qquad I\subseteq B\subseteq Y,
 \qquad\text{or}\qquad Y\subseteq B.
\end{equation}

\begin{proposition}\label{prop:basis-family}
There is a nonempty family $\mathcal B$ of admissible subsets of $E_0$
with the following properties:
\begin{enumerate}
 \item $\mathcal B$ is an antichain and is closed under complements.
 \item If $B\in\mathcal B$, $F\subseteq B$ and $H\subseteq E_0\setminus B$
 are finite with $|F|=|H|$, then $(B\setminus F)\cup H\in\mathcal B$.
 \item For every $I\subseteq Y\subseteq E_0$ with $Y\setminus I$ infinite,
 the alternative~\eqref{eq:interval-alternative} holds.
\end{enumerate}
\end{proposition}

\begin{proof}
As $E_0$ is countable, list all intervals with infinite gap in a sequence
of length $\mathfrak c$.
Recursively retain complementary pairs of admissible prototypes, requiring
the two directed differences of any distinct prototypes to be
$\mathcal J_0$-positive.
At any stage, call bases in the already retained local families
\emph{old bases}.
They are closed under complements.

Consider the current interval $I\subseteq Y$.
If~\eqref{eq:interval-alternative} already holds for an old basis,
skip this stage. Otherwise Proposition~\ref{prop:local-intervals}
implies that every old prototype $T$ satisfies
\begin{equation}\label{eq:unmet-interval}
 T\setminus I\notin\mathcal J_0,
 \qquad Y\setminus T\notin\mathcal J_0.
\end{equation}
Thus an unmet interval has the positive differences needed to reserve
separation from every old class.

An admissible set cannot contain an Upper set or be contained in a Lower
set, by Lemma~\ref{lem:zones}. Thus if $I$ is Upper, we must seek the
allowed basis below $I$; if $Y$ is Lower, we must seek it above $Y$.
Otherwise we seek a basis between the endpoints. Accordingly choose
\[
 [A,Z]=
 \begin{cases}
 [\varnothing,I],& I\text{ is Upper},\\
 [Y,E_0],&Y\text{ is Lower},\\
 [I,Y],&\text{otherwise}.
 \end{cases}
\]
The first two cases cannot coincide, by monotonicity of the Upper zone.
In every case $A$ is not Upper and $Z$ is not Lower.
We verify~\eqref{eq:available-separation} for all old prototypes.
In the third case it is~\eqref{eq:unmet-interval}.
In the first case $T\setminus A=T$ is positive.
If $I\setminus T$ were $\mathcal J_0$-small, $I\cap T$ would remain Upper and be
contained in $T$, forcing the admissible $T$ to be Upper.
Thus $Z\setminus T=I\setminus T$ is positive.
In the second case $Z\setminus T=T^c$ is positive.
If $T\setminus Y$ were $\mathcal J_0$-small, $T\cap Y$ would remain admissible and be
contained in the Lower set $Y$, a contradiction.
Thus $T\setminus A=T\setminus Y$ is positive as well.

At a stage $\alpha<\mathfrak c$ there are at most two prototypes per
earlier stage, hence fewer than $\mathfrak c$ old prototypes.
Lemma~\ref{lem:insert-prototype} therefore supplies an admissible $B$
between $A$ and $Z$ with both positive differences from every old
prototype.
Insert $B$ with the minus rule and $B^c$ with the plus rule.
The complement has both positive differences from every old $T$, because
$T^c$ is also old and
\[
 B^c\setminus T=T^c\setminus B,\qquad
 T\setminus B^c=B\setminus T^c.
\]
The two new prototypes have positive differences from one another by
Lemma~\ref{lem:zones}.
They therefore define new classes.
All probes of $B$ relative to itself vanish, so $B\in\mathcal B^-(B)$.
Its placement in $[A,Z]$ supplies one of the three outcomes
in~\eqref{eq:interval-alternative} for the original interval.

This recursion uses no regularity assumption on $\mathfrak c$.
For each individual $\alpha$ below its initial ordinal, the numbers of
predecessors, prototypes, and requested traces are all bounded by a fixed
cardinal less than $\mathfrak c$.
The reservation Lemma applies with that bound.
Choices can be made using a fixed well-order of $2^{E_0}$.

Let $\mathcal B$ be the union of all retained local basis families.
It is nonempty, since the interval $[\varnothing,E_0]$ receives an outcome.
Every member remains admissible under its $\mathcal J_0$-small change
from its prototype, and the family is closed under complements.
Lemma~\ref{lem:local-antichain} gives the antichain and finite-change
properties within each class.
For different classes, the positive directed differences of their
prototypes remain positive after $\mathcal J_0$-small changes.
Thus their bases are incomparable as well.
Finally every interval outcome persists, because no local family is
subsequently altered.
\end{proof}

\subsection{The maximal-extension axiom}

The interval construction is complete.
We now check that it produces a matroid, including maximal extension
inside arbitrary subsets; finite exchange alone does not supply this
last property.

\begin{theorem}\label{thm:local-matroid}
There is a matroid $Q$ on $E_0=L\dot\cup R$ whose bases are precisely
the family $\mathcal B$ in Proposition~\ref{prop:basis-family}.
It is equal to its dual on the same labels, and all its bases are admissible.
In particular, for every basis $B$,
\begin{equation}\label{eq:basis-rank-decrease}
 B_L\in\mathcal U
 \quad\Longrightarrow\quad
 D\setminus B_R\in\mathcal U
 \quad\text{and}\quad
 \rho(B_L)>\rho(D\setminus B_R).
\end{equation}
Every independent set of $Q$ that is not a basis can be extended by any
element outside it.
\end{theorem}

\begin{proof}
Define the independent sets to be all subsets of members of $\mathcal B$.
Since $\mathcal B$ is nonempty, this gives (I1), and heredity (I2) is
immediate.
The maximal independent sets are exactly $\mathcal B$:
each independent set is contained in a member, and the antichain property
prevents any such member from having a proper independent extension.

Suppose $I$ is independent but not maximal, and choose $B_0\in\mathcal B$
properly containing it.
Let $e\in E_0\setminus I$.
If $e\in B_0$, then $I\cup\{e\}$ is independent.
Otherwise choose $f\in B_0\setminus I$; the balanced swap
$(B_0\setminus\{f\})\cup\{e\}$ is a basis containing $I\cup\{e\}$.
This proves the final assertion of the theorem.
For any maximal independent $B_1$, its difference $B_1\setminus I$ is
nonempty, since $B_1\subseteq I$ would contradict its maximality.
Taking $e$ from this difference proves (I3).

For (IM), fix an independent $I\subseteq Y\subseteq E_0$.
If $Y\setminus I$ is finite, choose an independent extension of $I$
containing a maximum number of elements of that finite gap.
It is maximal within $Y$.
If the gap is infinite, use~\eqref{eq:interval-alternative}.
A basis $B\subseteq I$ forces $B=I$: put $I$ inside a basis $B'$ and
apply the antichain property to $B\subseteq I\subseteq B'$.
Thus in this case $I$ itself is maximal within $Y$.
A basis between $I$ and $Y$ is the required maximal extension.
A basis containing $Y$ makes $Y$ independent, so $Y$ itself is the
required extension.
This proves all four independent-set axioms.

The bases of the dual matroid are the complements of the bases of $Q$.
Since $\mathcal B$ is complement-closed, $Q^*=Q$ on $E_0$.
Admissibility was established in Proposition~\ref{prop:basis-family},
and~\eqref{eq:basis-rank-decrease} is its first case.
\end{proof}

\begin{remark}\label{rem:uniform-local}
The matroid $Q$ is uniform in the sense that every subset is independent
or spanning~\cite{BowlerGeschke2016}.
Indeed, if $X$ is dependent, maximize an independent subset $I$ inside
$X$ using (IM).
If $I$ were not a basis of $Q$, the last assertion of
Theorem~\ref{thm:local-matroid} would allow another element of $X$ to be
added.
Thus $X$ contains a basis and is spanning.
Every basis and its complement are infinite by admissibility, so $Q$
has infinite rank and corank.
Every finite $F\subseteq E_0$ is independent: starting from any basis,
insert the elements of $F$ that are missing and delete equally many
elements outside $F$, using its infinitude and balanced finite changes.
Yet $E_0$ is dependent, because every basis has a nonempty complement.
Thus $Q$ is not finitary, and self-duality shows that it is not cofinitary
either.
\end{remark}
